\chapter{Body Odor (Bromhidrosis)}

\section*{Definition}
Unpleasant body odor caused by bacterial breakdown of sweat and sebaceous secretions on the skin; apocrine bromhidrosis involves the axillae and genital region, while eccrine bromhidrosis affects the entire body.

\section*{When to See a Doctor}

\begin{itemize}
 \item Body odor changes suddenly without identifiable dietary or hygiene cause
 \item Fishy odor suggesting trimethylaminuria (genetic or acquired)
 \item Sweet, fruity breath with thirst, frequent urination, vomiting, abdominal pain, deep breathing, drowsiness, or confusion; seek emergency care because diabetic ketoacidosis is possible
 \item A new unusual odor with known liver or kidney disease, or with severe illness
 \item Accompanied by systemic symptoms (fever, weight loss, fatigue, hyperhidrosis)
\end{itemize}

\section*{How It Develops}
Sweat is usually odorless when produced. Odor develops when sweat and skin secretions are broken down by microorganisms on the skin. It is influenced by sweating, skin folds, hormones, clothing, medicines, and diet; a generalized unusual odor may occasionally reflect an underlying illness.

\section*{Causes}
\begin{itemize}
 \item Apocrine (axillary, genital): Normal variant with pungent odor, trimethylaminuria (fish-odor syndrome due to FMO3 mutation), hirsutism
 \item Eccrine (generalized): Obesity (intertriginous maceration), hyperhidrosis, diabetes mellitus (ketone body odor), hepatic failure (fetor hepaticus), renal failure (uremic fetor), diabetic ketoacidosis (fruity acetone odor)
 \item Dietary: Strong-smelling foods and alcohol can temporarily change odor
 \item Medicines or medical conditions that increase sweating or alter metabolism
\end{itemize}

\section*{Related Symptoms}
\begin{itemize}
 \item Hyperhidrosis (excessive sweating)
 \item Intertrigo from maceration
 \item Skin maceration or fungal infection in skin folds
 \item Halitosis (bad breath)
 \item Social withdrawal or anxiety
\end{itemize}
\textbf{Important:} Trimethylaminuria is a rare genetic disorder causing fish-odor syndrome that can be managed with dietary restriction of choline-rich foods.

\section*{Medical Evaluation}
\begin{itemize}
 \item Clinical history and physical exam distinguishing focal from generalized body odor.
 \item Testing is guided by the history and examination. Trimethylamine testing is considered only when a persistent fish-like odor affects sweat, breath, and urine and common causes have been excluded.
 \item Blood glucose and ketones are checked urgently when diabetic ketoacidosis is suspected; liver or kidney tests are considered when symptoms suggest those conditions.
 \item A skin swab is used selectively when a bacterial or fungal infection is suspected.
\end{itemize}

\section*{Treatment Options}
\begin{itemize}
 \item Stronger antiperspirants may be prescribed for excessive sweating. Topical antibiotics are not routine and are used only for a diagnosed skin infection.
 \item If severe hyperhidrosis persists, a clinician may consider medicines, botulinum toxin, or other specialist treatments; benefits and side effects vary.
 \item Laser hair removal and microwave thermolysis for sweat gland destruction in refractory cases.
 \item Surgical: Subcutaneous curettage or liposuction-curettage of apocrine glands; surgical excision reserved for severe unresponsive cases.
\end{itemize}

\section*{Self-Care and Home Management}
\begin{itemize}
 \item Wash the armpits, groin, and feet regularly with a gentle cleanser and dry skin folds thoroughly; harsh antibacterial products can irritate skin.
 \item Use antiperspirants (aluminum chloride) to reduce sweat production, not just deodorants which mask odor.
 \item Wear breathable natural fibers (cotton, wool) and change clothing daily.
 \item Wear clean, loose, breathable clothing and change it regularly. Hair removal is optional and may irritate the skin.
 \item Avoid dietary triggers: reduce garlic, onions, curry, alcohol, and cruciferous vegetables if sensitive.
\end{itemize}

\section*{References}
\begin{thebibliography}{99}
\bibitem{body_odor:ref1} Kanlayavattanakul M, Lourith N. ``Body malodours and their topical treatment agents.'' \textit{Int J Cosmet Sci}. 2011;33(4):298-311.
\bibitem{body_odor:ref2} James AG, Austin CJ, et al. ``The biological basis of malodour.'' \textit{J Appl Microbiol}. 2013;115(6):1261-1272.
\bibitem{body_odor:ref3} Guillet G, Zampetti A, et al. ``Bromhidrosis: a review.'' \textit{J Eur Acad Dermatol Venereol}. 2007;21(5):598-604.
\bibitem{body_odor:ref4} James WD, Elston DM, et al. \textit{Andrews\' Diseases of the Skin: Clinical Dermatology}, 13th ed. Elsevier, 2019.
\bibitem{body_odor:ref5} National Health Service. ``Body odour.'' Reviewed May 2025. https://www.nhs.uk/symptoms/body-odour-bo/.
\bibitem{body_odor:ref6} MedlinePlus Genetics. ``Trimethylaminuria.'' Accessed September 2026. https://medlineplus.gov/genetics/condition/trimethylaminuria/.
\end{thebibliography}
